\section{Introduction}\label{ch:introduction}

Australian energy prices are unremarkable, both in our minds and in scale. Reserved for the eyes of energy retailers, the wholesale price averaged less than $\$100$ in all states that the National Energy Market serviced in the last quarter. \cite{aercharts} Within short time frames, the regional price can jump by one to two orders of magnitude, before it is hard capped by the market regulator. Such extrema are the driving force behind the market's economic actors - retailers who seek to insure themselves against paying such high prices, and generators who will vary generation on \emph{price spike} forecasts. The nature of this report is to answer such: \emph{Given everything observable at the end of a dispatch interval, how well can we predict that the price will be above \$300/MWh some moments later?} 

\subsection{The National Electricity Market}\label{sec:NEM}

\subsection{Price Spikes}\label{sec:pricespike}
Why do we define the threshold to be \$300/MWh? In some ways it is arbitrary - it is no longer the administered price cap \cite{aemc2026}, and it doesn't compare to the most extreme prices that are reached. However, it was chosen because the threshold is crossed often enough, and also because it is the price level of the standard cap contracts - the call option retail energy providers purchase to insure against risk. \cite{aemcmarkets} 

\subsection{Aim}\label{sec:aim}
